\section{How the Analysis Works}

The method follows one rule from beginning to end: establish the object first, perform the arithmetic second, and interpret it third. This order matters because an attractive result can vanish when the underlying Hebrew form is corrected.

\subsection{Start with a Definite Object}

Each entry must have a stated consonantal form, a textual context, and an object type. Before any number is calculated, the analysis fixes articles, construct spelling, word division, and final-letter position. Punctuation, cantillation, and vowel points receive no value.

The inherited manuscript contained 43 derivation objects, but they were not all the same kind of thing. The source audit separates them into:

\begin{itemize}
\item 24 primary biblical names and titles;
\item 5 later traditional forms;
\item 4 messianic or Christian-comparative forms;
\item 4 foreign, constructed, or false-match rows that are excluded from name counts;
\item 6 Kabbalistic formulas or expansion labels.
\end{itemize}

The phrase \emph{Christian-comparative} identifies the premise by which an entry enters the analysis. Within a Christian reading, Yeshua is included as a divine personal name; the separate label keeps the Hebrew-Bible-only comparison reproducible.

I retained the original 43-object boundary in the derivation appendix because I did not want correction to erase the history of the study. A separate registry then tested 82 additional names, titles, and complete expressions. Every candidate remains in that registry whether or not it produced a pattern. This keeps later discoveries from quietly changing the primary set.

\subsection{Apply the Same Sequence}

Every admitted literal form passes through the same sequence:

\begin{enumerate}
\item calculate standard gematria;
\item calculate Mispar Gadol when a genuine final letter is present;
\item convert both values to base 12;
\item record the decimal and base-12 digit-sum paths;
\item test decimal and base-12 palindrome status;
\item note exact shared values, canonical contacts, and exploratory mathematical echoes.
\end{enumerate}

The basis is carried through every line. A visible 72 in base 12 is not silently treated as decimal 72, and a formula label is not counted as though it were an independently discovered divine name.

A later exploratory audit added one operation that was not part of the original sequence. For an odd-length base-12 palindrome, the middle digit is removed and the remaining ordered digits are compared with the complete base-12 values of other admitted expressions. The audit calls this a \emph{named outer frame}. Standard gematria and Mispar Gadol are checked separately before any mixed-layer comparison is considered. A complete follow-up then enumerates all 132 forms $ada_{12}$ with $a\neq0$ and $d$ any base-12 digit. Because the operation was motivated by known results, it can establish exact finite structure but not a discovery probability.

\subsection{Let Corrections Count Against the Pattern}

The source rule is useful only if it can reject a result I would have preferred to keep. Several corrections do exactly that. Five inherited Elohei spellings contained a vav that was absent from the cited biblical forms. The supposed El HaShamayim and El HaAretz titles came from passages where \emph{el} functions as the preposition ``to.'' The former equality between El Olam and El HaNe'eman depended on dropping the article found in Deuteronomy 7:9.

Those losses are part of the evidence. At the same time, the audit preserved or clarified strong results such as YHWH, Shaddai, Ehyeh Asher Ehyeh, El Olam, Yeshua, and Yeshua HaMashiach. A method that keeps only its successes would make the paper easier to write and harder to trust.

\subsection{Keep Observation and Meaning in View}

For each result, I distinguish the following levels:

\begin{description}[leftmargin=3.3cm]
\item[Calculated] The Hebrew value, base conversion, digit sum, or palindrome can be reproduced.
\item[Classified] The calculation belongs to a stated layer, such as terminal 8 or base-12 palindrome.
\item[Interpreted] A structural or theological meaning is proposed for the result.
\item[Exploratory] A promising relationship has been noticed but was not part of the frozen primary test.
\item[Unresolved] The form, source, history, or interpretation still needs review.
\end{description}

These labels are not meant to drain the paper of wonder. They allow a striking interpretation to be stated without disguising it as arithmetic.

\subsection{Compare the Pattern with Other Hebrew Phrases}

The selected names were also compared with structurally similar Hebrew Bible phrases. Each target was matched by word count, word-length pattern, and number of final letters. One hundred thousand simulated corpora were then tested for palindromes across bases 2 through 20 under standard gematria and Mispar Gadol.

This comparison asks whether phrases with the same broad shape produce as many palindromes. The results are explained near the end of the main paper; the exact protocol, probabilities, control pool, and reproducibility details appear in Appendix D.

The controls deliberately replace names and declarations with other structures. That makes them useful for testing arithmetic recurrence, but it also limits what they can imitate. An ordinary phrase with the same value is not thereby the same object as a divine name, and a favorable alternative anchor set is not thereby a set with the same canonical associations. Throughout the interpretation, the complete attested name, title, or declaration remains the primary object. Splitting it into components or substituting a synthetic form such as $X$--Asher--$X$ is a diagnostic test of the mechanism, not a replacement for the expression being studied.

\subsection{Leave an Audit Trail}

The current Hebrew forms live in a source-controlled registry, and the derivations are generated from the corrected results table. Biblical calculations use a locked Open Scriptures Hebrew Bible revision, while representative passages were checked through Sefaria.\cite{oshb,sefaria_library} The inherited claims are preserved separately so that a reader can see what changed. Source corrections appear with the full rows in Appendix A; the expansion findings and formal layer counts follow in Appendices B and C. The complete machine-readable registries remain in the repository.

This architecture gives the paper two reading levels. A reader can follow the argument in the main text without learning the repository machinery. A skeptical or technical reader can inspect every operation afterward. Both readers are looking at the same calculations.
